\section{Introduction}
\label{sec:intro}

\noindent
What the robot fails to see, it fails to do.
Modern robot learning places large pretrained perception models at this first point of contact with the physical world.
Vision-Language-Action (VLA) models ($\pi_0$~\citep{black2024pi0}, OpenVLA~\citep{kim2024openvla}, GR00T~\citep{bjorck2025gr00t}) inherit visual encoders from Vision-Language Models (VLMs) such as SigLIP~\citep{zhai2023siglip} and DINOv2~\citep{oquab2023dinov2}.
Modular systems chain open-vocabulary detectors (GroundingDINO~\citep{liu2024groundingDINO}), segmenters (SAM2~\citep{ravi2024sam2}), and depth models (ZoeDepth~\citep{bhat2023zoedepth}) for grasp planning~\citep{liu2024okrobot, ze2024dp3}.
World models (TesserAct~\citep{zhen2025tesseract}) consume and generate depth and normal maps conditioned on actions.
In each case, perception quality strongly influences downstream success.

\paragraph{The unjustified backbone choice.}
Despite this dependence, backbone selection is ad hoc.
Li~et~al.~\citep{li2025robovlms} confirmed that visual backbone choice dominates VLA policy success---more than architecture or training recipe---but evaluated entirely in simulation.
OK-Robot~\citep{liu2024okrobot} showed that perception quality is the dominant bottleneck in real-world pick-and-place (success drops from 82\% to 59\% under clutter), but evaluated only one perception pipeline without comparing alternatives.
The problem is not a lack of models---it is a lack of evidence about which model to use.
Three deficiencies in current evaluation keep the choice unjustified.
\textbf{(i)~Sensor mismatch:} models are ranked on COCO~\citep{lin2014coco} and ImageNet~\citep{deng2009imagenet}---curated internet images at high resolution. Robots see through \$30 Intel RealSense D435 sensors at $640\!\times\!480$ with structured-light depth noise, motion blur, and specular artifacts. No benchmark evaluates models on this imagery.
\textbf{(ii)~Task silos:} each task is evaluated on a different benchmark with different scenes, so per-task leaderboards cannot compose into a deployment recommendation.
\textbf{(iii)~No scene-state variation:} existing benchmarks capture static scenes. A robot's scene changes as it acts---objects move, hands occlude, clutter shifts. No benchmark measures how perception degrades during manipulation and whether it recovers afterward.

\paragraph{Evaluation methodology, not just a dataset.}
The 2026 NeurIPS E\&D track asks: \textit{what claims does an evaluation support, under what assumptions?}
We argue that answering ``which perception backbone should a robot use?'' requires \textbf{evaluation methodology} that existing benchmarks lack---not more data, but better measurement.
Specifically, evaluation must:
(a)~create controlled scene-state variation so that performance differences are primarily attributable to scene state, not confounded by object or lighting changes;
(b)~measure properties beyond accuracy---temporal consistency, phase-transition robustness, geometric coherence---that determine whether a model \textit{works in deployment};
and (c)~stratify difficulty by perceptual ambiguity, not post-hoc geometry.

\paragraph{\coolname{}: methodology + dataset + toolkit.}
We introduce \coolname{}, which provides all three.
The \textit{three-phase capture protocol} (\S\ref{sec:three-phase}) mirrors the robot manipulation cycle (observe $\to$ act $\to$ verify), holding scene identity fixed while varying scene state.
\textit{Effort-Stratified Difficulty} (\S\ref{sec:features}) derives difficulty from annotation effort---a direct signal of perceptual ambiguity.
Three \textit{deployment-readiness metrics} (\S\ref{sec:deploy-metrics})---Temporal Stability (TS), State-Transition Robustness (STR), and Geometric Coherence (SGC)---compose into an Embodied Deployment Score (EDS, \S\ref{sec:eds}) that ranks models for real-world use.
These are instantiated on a 75{,}000-frame real-world RGB-D corpus across 100 scenes and ${\sim}$70 object categories (\S\ref{sec:dataset}).
All five tasks---depth, detection/grounding, tracking, scene QA, spatial QA---are evaluated on the \textit{same} scenes and objects, described in crowd-sourced natural language~\citep{padalunkal2023fewsol}, with a pip-installable toolkit (\S\ref{sec:toolkit}) that enables reproducible bring-your-own-model evaluation.

\paragraph{Human proxy for robot perception.}
A human carries the exact sensor rig a mobile robot would use---D435 (RGB-D at $640\!\times\!480$) + T265 (stereo tracking) + GoPro (egocentric)---through 100 real-world scenes: kitchens, offices, outdoor paths, staircases, day and night, reflective and transparent surfaces.
This is not a limitation but a deliberate design: the learning-from-humans paradigm~\citep{padalunkal2024iteach, 2025hrt1, khazatsky2024droid, kareer2024egomimic} already uses human-captured video as robot training data.
If pretrained models fail on human-captured RealSense imagery, they fail on the data robots are trained on.
Paired single-object and multi-object scenes enable controlled comparison of text-prompted vs.\ image-prompted perception under real-world conditions.

% ── TODO(results): replace once experiments complete ──
\paragraph{Key findings.}
Evaluating \todo{N} models across five tasks, we find:
\todo{REPLACE with 2--3 concrete findings. Candidates:
(i) static accuracy is weakly correlated with deployment robustness ($\rho = X$);
(ii) model rankings on Hard scenes diverge from Easy ($\tau = X$);
(iii) VLM spatial QA accuracy drops by X\% during interaction, directly implicating VLA backbone robustness;
(iv) EDS predicts downstream grasp success ($\rho = X$).}

\paragraph{Contributions.}
\begin{enumerate}[leftmargin=*, topsep=2pt, itemsep=1pt]
    \item \textbf{An evaluation methodology for deployment-ready robot perception.}
    Three-phase protocol + ESD + task-specific deployment-readiness metrics (TS, STR, SGC) + EDS---a methodological framework that can be applied to future datasets and tasks beyond \coolname{}.

    \item \textbf{The RPX dataset and multi-task benchmark.}
    75{,}000-frame real-world RGB-D corpus across 100 scenes and ${\sim}$70 object categories, with five perception tasks (depth, detection/grounding, tracking, scene QA, spatial QA) evaluated on the \textit{same} scenes and objects---enabling the first cross-task deployment recommendations on shared real-world scenes.

    \item \textbf{Empirical findings.}
    \todo{2--3 sentence findings summary, including the named finding.}

    \item \textbf{Self-sustaining community toolkit.}
    \texttt{rpx-benchmark}: pip-installable, bring-your-own-model, with plugin registries for new tasks and metrics. 138 CI-verified tests.
\end{enumerate}

\paragraph{Scope.}
\coolname{} tests the claim that pretrained perception backbones differ in deployment robustness in ways not predictable from standard accuracy alone.
It evaluates perception at the output level (depth maps, boxes, QA answers); for VLAs that consume encoder features, output-level metrics serve as proxies for feature quality (justified in \S\ref{sec:eds}, validated in \S\ref{sec:downstream}).
Assumptions and scope limitations are detailed in Appendix~\ref{sec:appendix-scope}.
